\documentclass[conference]{IEEEtran}
\IEEEoverridecommandlockouts

\usepackage[utf8]{inputenc}
\usepackage[T1]{fontenc}
\usepackage[english]{babel}
\usepackage{amsmath,amssymb,amsfonts}
\usepackage{graphicx}
\usepackage{booktabs}
\usepackage{array}
\usepackage{cite}
\usepackage{xcolor}
\usepackage[table]{xcolor}
\usepackage{url}
\usepackage{algorithm}
\usepackage{algorithmicx}
\usepackage{algpseudocode}
\usepackage{pgfplots}
\pgfplotsset{compat=1.18}
\usepackage{tikz}
\usetikzlibrary{shapes,arrows,positioning,calc}
\usepackage[hidelinks]{hyperref}

\begin{document}
	
	\title{Dynamic Parameter Adaptation in LR-FHSS LoRaWAN Networks\\Using Deep Reinforcement Learning}
	
	\author{
		\IEEEauthorblockN{Marthy Elys\'ee BAZOUM}
		\IEEEauthorblockA{Universit\'e Nazi BONI (UNB), \\ Bobo-Dioulasso, Burkina Faso}
		\and
		\IEEEauthorblockN{Tiguiane YELEMOU}
		\IEEEauthorblockA{Universit\'e Nazi BONI (UNB), \\ Bobo-Dioulasso, Burkina Faso}
		\and
		\IEEEauthorblockN{G\'erard CHALHOUB}
		\IEEEauthorblockA{\textit{Universit\'e Clermont Auvergne}\\
			Clermont-Ferrand, France}
		\and
		\IEEEauthorblockN{Hamadoun TALL}
		\IEEEauthorblockA{Universit\'e Nazi BONI (UNB),\\ Bobo-Dioulasso, Burkina Faso}
	}
	
	\maketitle
	
	% =====================
	% ABSTRACT
	% =====================
	\begin{abstract}
		The rapid expansion of the Internet of Things (IoT) requires Low-Power Wide
		Area Networks (LPWAN) to meet growing demands for capacity and reliability.
		LR-FHSS (Long Range -- Frequency Hopping Spread Spectrum), recently introduced
		to extend LoRaWAN, enables massive and resilient communications, particularly
		for direct satellite links. However, dynamically optimizing its transmission
		parameters (Data Rate and transmit power) under varying radio channels remains
		a major open challenge, as existing static or heuristic approaches are
		unsuitable. This paper proposes an intelligent adaptation strategy based on a
		Double Deep Q-Network (DDQN) agent to simultaneously maximize reliability and
		energy efficiency of LR-FHSS transmissions. Our approach comprises three
		contributions: (i) a high-fidelity event-driven Python simulator, rigorously
		validated against real-world measurements (RMSE of 12.14\% for PDR and
		6.56\,dB for RSSI); (ii) a formal Markov Decision Process (MDP) problem
		statement with a 76-action discrete space (DR $\times$ power); and (iii) a
		trained DDQN agent evaluated against static configurations. Results show that
		the agent improves the average PDR by +2.48 percentage points, with gains
		reaching +15.6\% at 4000\,m, while reducing average transmit power by 32\%
		and energy consumption per packet by 4.9\% on average (up to 62.5\% at short
		range).
	\end{abstract}
	
	\begin{IEEEkeywords}
		LR-FHSS, LoRaWAN, IoT, Deep Reinforcement Learning, DDQN, Adaptive
		Transmission, Energy Efficiency, LPWAN
	\end{IEEEkeywords}
	
	% =====================
	% I. INTRODUCTION
	% =====================
	\section{Introduction}
	
	The massive expansion of the Internet of Things (IoT) places increasing demands
	on Low-Power Wide Area Networks (LPWANs) in terms of capacity and reliability.
	While LoRaWAN has established itself as a reference technology thanks to its
	Chirp Spread Spectrum (CSS) modulation~\cite{Raza2017LPWANSurvey}, its
	limitations in scalability and interference resistance motivated the development
	of LR-FHSS (Long Range -- Frequency Hopping Spread Spectrum). Standardized by
	the LoRa Alliance in 2020~\cite{loraalliance2020regional}, LR-FHSS replaces
	CSS with GMSK modulation combined with intra-packet frequency hopping. It
	fragments messages over ultra-narrow 488\,Hz sub-channels, promises a tenfold
	increase in network capacity, and is particularly suited for direct-to-satellite
	IoT links~\cite{ieee_lr_fhss_2021}.
	
	Despite these advances, the dynamic optimization of LR-FHSS transmission
	parameters (Data Rate and transmit power) under varying channel conditions
	remains a major open problem. Existing adaptive approaches, such as
	ADR~\cite{Li2018ADRanalysis} or RL-based solutions like
	LoRa-DDPG~\cite{greenlora,greenlora2024}, optimize the Spreading Factor
	(SF)---a parameter specific to CSS modulation that does not exist in
	LR-FHSS---and rely on the LoRaSim physical model, which does not capture
	LR-FHSS fragment-level collision mechanics.
	
	This paper addresses this gap with three main contributions:
	\begin{itemize}
		\item A \textbf{high-fidelity event-driven simulator} that faithfully
		implements LR-FHSS (RP002-1.0.5~\cite{loraalliance2025regional}), including
		fragment-level collision detection and a realistic propagation model,
		validated against experimental data~\cite{lrfhss_comparaison};
		\item A \textbf{formal MDP problem formulation} for joint DR and transmit
		power optimization, with a multi-objective reward function balancing
		reliability, energy efficiency, and spectral occupancy;
		\item A \textbf{DDQN agent} trained on the simulator and evaluated against
		static per-DR configurations across 13 distances from 139\,m to 4000\,m.
	\end{itemize}
	
	The remainder of this paper is structured as follows.
	Section~\ref{sec:background} presents LR-FHSS technical foundations.
	Section~\ref{sec:related} reviews related work.
	Section~\ref{sec:system} describes the MDP formulation and DDQN design.
	Section~\ref{sec:eval} presents results.
	Section~\ref{sec:conclusion} concludes.
	
	% =====================
	% II. LR-FHSS BACKGROUND
	% =====================
	\section{LR-FHSS Technical Background}
	\label{sec:background}
	
	\subsection{Modulation and Architecture}
	
	LR-FHSS represents a fundamental departure from LoRa CSS. Instead of
	adjusting the Spreading Factor for link robustness, it uses GMSK modulation
	and achieves frequency diversity through intra-packet frequency hopping. Each
	data fragment occupies a bandwidth of 488\,Hz (Occupied Bandwidth, OBW), and a
	complete packet transmission hops across sub-channels within an Operating
	Channel Width (OCW) of 137\,kHz or 336\,kHz~\cite{semtech2022lrfhss}. This
	approach reduces collision probability in dense networks and improves resilience
	to narrowband interference.
	
	\subsection{Data Rate Configurations}
	
	For the EU868 band, the LoRa Alliance specification~\cite{loraalliance2025regional}
	defines four LR-FHSS Data Rates (DR8--DR11), detailed in Table~\ref{tab:dr_config}.
	Configurations with Coding Rate CR~1/3 (DR8, DR10) prioritize robustness at the
	cost of throughput, while CR~2/3 (DR9, DR11) doubles the useful bit rate with
	less error protection. The wider OCW of DR10/DR11 (336\,kHz) provides 688
	available sub-channels vs. 280 for DR8/DR9, improving anti-collision capacity.
	
	\begin{table}[!t]
		\centering
		\caption{LR-FHSS Data Rate Configurations (EU868)}
		\label{tab:dr_config}
		\begin{tabular}{|c|c|c|c|c|}
			\hline
			\textbf{DR} & \textbf{OCW (kHz)} & \textbf{CR} & \textbf{Channels} & \textbf{Rate (bps)} \\
			\hline
			DR8  & 137 & 1/3 & 280 & 162 \\
			DR9  & 137 & 2/3 & 280 & 325 \\
			DR10 & 336 & 1/3 & 688 & 162 \\
			DR11 & 336 & 2/3 & 688 & 325 \\
			\hline
		\end{tabular}
	\end{table}
	
	\subsection{Time-on-Air}
	
	The total Time-on-Air (ToA) for a packet with payload length $L$ bytes is:
	\begin{equation}
		T_{\text{air}} = N_h \times 233.472\,\text{ms}
		+ \left\lceil \frac{N_{\text{bits}}}{48} \right\rceil \times 102.4\,\text{ms}
		\label{eq:toa}
	\end{equation}
	where $N_h = 3$ for CR 1/3 and $N_h = 2$ for CR 2/3, and $N_{\text{bits}}$ is
	the total number of coded bits after convolution coding.
	
	% =====================
	% III. RELATED WORK
	% =====================
	\section{Related Work}
	\label{sec:related}
	
	\textbf{Receiver-level optimization.}
	Maldonado et al.~\cite{maldonado2025headerless} propose a receiver capable of
	decoding LR-FHSS frames with all headers lost, achieving a 50\% gain in
	decoded packet rate. This requires significant gateway hardware modification
	and does not address dynamic parameter adaptation.
	
	\textbf{FHS and demodulator optimization.}
	The same group~\cite{maldonado2025enhancing} reverse-engineered commercial
	LR-FHSS hopping sequences (LFSR-based m-sequences~\cite{lfsr}) and proposed
	optimized families (Li-Fan, Lempel-Greenberger), yielding up to 100\% capacity
	gain combined with Early-Decode/Early-Drop demodulator allocation. These
	improvements are static (designed offline) and channel-agnostic.
	
	\textbf{Doppler compensation.}
	Ullah et al.~\cite{ullah_extending_doppler_2025} present beacon-based Doppler
	pre-compensation for satellite links, eliminating frequency-induced losses with
	only 2.5\% energy overhead. This addresses a specific impairment but provides no
	DR or power adaptation mechanism.
	
	\textbf{RL-based optimization for LoRa.}
	Saleem et al.~\cite{greenlora,greenlora2024} apply DDPG to jointly optimize
	(Channel, SF, TX Power, Time Slot) in LoRaWAN, reporting ×98 energy efficiency
	gains. However, these works confuse LoRa and LR-FHSS: the optimized SF parameter
	does not exist in LR-FHSS, and the CSS physical model is incompatible with
	LR-FHSS fragment-level mechanics. Their results cannot be transposed to LR-FHSS.
	
	\textbf{Positioning.}
	To the best of our knowledge, no existing work proposes a dynamic, multi-objective
	optimization strategy specifically designed for the physical parameters of LR-FHSS.
	Our work fills this gap by combining a validated LR-FHSS simulator with a DDQN
	agent operating over LR-FHSS-compatible (DR, power) action pairs.
	
	% =====================
	% IV. SYSTEM DESIGN
	% =====================
	\section{System Model and DDQN Design}
	\label{sec:system}
	
	\subsection{MDP Formulation}
	
	The LR-FHSS adaptation problem is formalized as a Markov Decision Process
	$(\mathcal{S}, \mathcal{A}, \mathcal{P}, \mathcal{R}, \gamma)$~\cite{puterman1994mdp,sutton2018reinforcement},
	where an agent located on an edge server learns a transmission policy through
	interaction with the network simulator.
	
	\textbf{State space $\mathcal{S}$:}
	The state $s_t \in \mathbb{R}^8$ encodes the channel and node context before
	each transmission:
	\begin{equation}
		s_t = \bigl[d_{\text{n}},\, \overline{\text{SNR}},\, \tau_s,\, r_n,\,
		\log(d)_n,\, t_d,\, \sigma_{\text{SNR}},\, L_n\bigr]
	\end{equation}
	where $d_{\text{n}}$ is the normalized distance (max 4\,km),
	$\overline{\text{SNR}}$ the moving-average SNR over the last 3 transmissions,
	$\tau_s$ the success rate over the last 5 attempts, $r_n$ the normalized
	retry count, $\log(d)_n$ the log-distance (linearizing path loss),
	$t_d$ a time-of-day factor, $\sigma_{\text{SNR}}$ link stability (inverse
	SNR std), and $L_n$ the normalized payload size. All components are bounded
	to $[0,1]$ or $[-1,1]$ for stable neural network training.
	
	\textbf{Action space $\mathcal{A}$:}
	Each action jointly selects a Data Rate and transmit power:
	\begin{equation}
		\mathcal{A} = \{\text{DR8},\ldots,\text{DR11}\}
		\times \{-4, -3, \ldots, 13, 14\}\,\text{dBm}
	\end{equation}
	yielding $|\mathcal{A}| = 4 \times 19 = 76$ discrete actions, encoded as
	unique integers $a \in \{0,\ldots,75\}$.
	
	\textbf{Reward function $\mathcal{R}$:}
	The reward trades off reliability against energy and spectral efficiency:
	\begin{equation}
		\mathcal{R} =
		\begin{cases}
			-200 & \text{(failure)} \\[4pt]
			100 + \mathcal{R}_{\text{ToA}} + \mathcal{R}_{\text{pwr}}
			+ \mathcal{R}_{\text{retry}} + \Omega & \text{(success)}
		\end{cases}
	\end{equation}
	The penalty $\phi = -200$ is twice the success base $\rho = +100$, forcing the
	agent to prioritize packet delivery. Sub-penalties are:
	\begin{align}
		\mathcal{R}_{\text{ToA}} &= -80 \cdot \min\!\left(\tfrac{T_{\text{air}}}{T_{\max}}, 1\right),
		\quad T_{\max} = 3813\,\text{ms} \\
		\mathcal{R}_{\text{pwr}} &= -50 \cdot \tfrac{P_{\text{TX}}+4}{18} \\
		\mathcal{R}_{\text{retry}} &= -30 \cdot \tfrac{N_{\text{retry}}}{5}
	\end{align}
	The term $\Omega$ adds a SNR quality bonus (encouraging the ``just enough''
	operating point, $-5 \leq \text{SNR} \leq 5$\,dB) and a DR efficiency bonus
	(favoring CR 2/3 when the channel permits). Weight analysis confirms that
	reliability accounts for 53.3\% of the total reward range, energy and spectral
	efficiency for 41.4\%.
	
	\subsection{DDQN Architecture}
	
	The DDQN algorithm~\cite{mnih2015dqn,vanhasselt2016ddqn} decouples action
	selection from evaluation to reduce overestimation bias present in vanilla DQN:
	\begin{equation}
		y_t = r_t + \gamma\,Q_{\theta^-}\!\!\left(s_{t+1},\,
		\arg\max_{a'} Q_\theta(s_{t+1}, a')\right)
	\end{equation}
	
	The neural network is a feedforward network with: input layer (8 neurons) $\to$
	FC-128 (ReLU) $\to$ Dropout ($p = 0.1$) $\to$ FC-64 (ReLU) $\to$ output (76
	linear neurons). The network has $\approx$14,000 parameters, enabling
	millisecond-level inference on standard hardware.
	
	\textbf{Training.}
	The agent is trained over 3,000 episodes of 100 steps each, with:
	experience replay buffer (capacity 10,000),
	batch size 64,
	learning rate $\alpha = 10^{-4}$,
	discount factor $\gamma = 0.99$,
	soft target update $\tau = 0.001$,
	$\epsilon$-greedy exploration decaying from 1.0 to 0.05 with factor 0.995.
	At each episode, the node distance is sampled uniformly in $[0, 4]$\,km,
	forcing the agent to learn a generalist policy. Training stabilizes around
	episode 450, reaching an average success rate of 89\% over the final 500
	episodes, with reward converging to $6870 \pm 2150$.
	
	% =====================
	% V. EVALUATION
	% =====================
	\section{Evaluation Results}
	\label{sec:eval}
	
	\subsection{Simulator Validation}
	
	The simulator was validated by comparing simulated PDR and RSSI against
	real LR-FHSS measurements from~\cite{lrfhss_comparaison,sanchez_vital_2025_data}
	for all four DR configurations. Simulations used $P_{\text{TX}} = 14$\,dBm, 10
	nodes, 30-minute runs ($\approx$400 packets per configuration), with propagation
	parameters $n = 3.3$ and $\sigma_s = 7.0$\,dB. Table~\ref{tab:validation}
	summarizes the RMSE results.
	
	\begin{table}[!t]
		\centering
		\caption{Simulator Validation: RMSE vs.\ Real Measurements}
		\label{tab:validation}
		\begin{tabular}{|c|c|c|}
			\hline
			\textbf{Data Rate} & \textbf{RMSE PDR (\%)} & \textbf{RMSE RSSI (dB)} \\
			\hline
			DR8  & 10.69 & 7.21 \\
			DR9  & 21.40 & 6.84 \\
			DR10 & 5.11  & 5.92 \\
			DR11 & 11.38 & 6.28 \\
			\hline
			\textbf{Average} & \textbf{12.14} & \textbf{6.56} \\
			\hline
		\end{tabular}
	\end{table}
	
	The simulator correctly reproduces the performance hierarchy (DR8/DR10 $>$
	DR9/DR11 at long range) and the characteristic PDR decay shape, validating
	its use as a training and evaluation environment.
	
	\subsection{PDR Performance}
	
	Table~\ref{tab:pdr} compares the PDR of static configurations (mean over
	DR8--DR11 at 14\,dBm) and the DDQN agent.
	
	\begin{table}[!t]
		\centering
		\caption{PDR Comparison: Static vs.\ DDQN Agent}
		\label{tab:pdr}
		\begin{tabular}{|c|c|c|c|}
			\hline
			\textbf{Dist. (m)} & \textbf{Static (\%)} & \textbf{DDQN (\%)} & \textbf{Gain (pts)} \\
			\hline
			139   & 99.9 & 100.0 & +0.1 \\
			331   & 99.7 & 99.0  & \cellcolor{red!20}$-$0.7 \\
			775   & 99.2 & 99.0  & \cellcolor{red!20}$-$0.2 \\
			1015  & 98.7 & 99.0  & +0.3 \\
			1194  & 97.1 & 97.7  & +0.6 \\
			1260  & 97.5 & 98.3  & +0.8 \\
			1753  & 95.7 & 99.2  & +3.5 \\
			1977  & 89.1 & 92.0  & +2.9 \\
			2000  & 97.1 & 97.2  & +0.1 \\
			2816  & 72.6 & 78.4  & +5.8 \\
			3000  & 87.7 & 93.2  & +5.5 \\
			4000  & 72.4 & 83.7  & \cellcolor{green!20}+11.3 \\
			\hline
			\textbf{Avg.} & \textbf{92.7} & \textbf{95.0} & \cellcolor{green!20}\textbf{+2.3} \\
			\hline
		\end{tabular}
	\end{table}
	
	The DDQN agent improves PDR in 10 out of 12 reported distances (83\%),
	with the largest gains at long range: +11.3\,pts at 4000\,m (+15.6\%
	relative), +5.8\,pts at 2816\,m, and +5.5\,pts at 3000\,m. Minor
	regressions at 331\,m and 775\,m ($\leq$0.7\,pt) are negligible given
	the already near-perfect PDR ($>$99\%) at those distances.
	
	\subsection{Adaptive Strategy Analysis}
	
	The DDQN's DR selection follows a clear distance-aware policy:
	at distances $\leq$331\,m, DR9/DR11 (CR 2/3, high throughput) are selected
	exclusively, yielding the maximum LR-FHSS throughput of 142.05\,bps. Beyond
	775\,m, the agent transitions entirely to DR8/DR10 (CR 1/3, robust), dropping
	to 89.59\,bps to preserve link reliability. This behavior is consistent with
	optimal communication theory: exploit high-efficiency configurations when the
	channel is favorable, and fall back to robust ones as conditions degrade.
	
	\subsection{Transmit Power and Energy Efficiency}
	
	The DDQN achieves a \textbf{32\% average transmit power reduction} (from
	14.0\,dBm to 9.5\,dBm). Power savings are particularly dramatic at short
	range: at 139\,m, the agent selects 3.0\,dBm (78.6\% reduction). Beyond
	2000\,m, power stabilizes at 12.0\,dBm, and reaches 14.0\,dBm only at 4000\,m.
	
	Table~\ref{tab:energy} reports energy consumption per packet.
	
	\begin{table}[!t]
		\centering
		\caption{Energy Per Packet: Static vs.\ DDQN Agent}
		\label{tab:energy}
		\begin{tabular}{|c|c|c|c|}
			\hline
			\textbf{Dist. (m)} & \textbf{Static (mJ)} & \textbf{DDQN (mJ)} & \textbf{Variation} \\
			\hline
			139  & 76.88 & 28.81 & \cellcolor{green!20}$-$62.5\% \\
			331  & 76.92 & 29.62 & \cellcolor{green!20}$-$61.5\% \\
			775  & 76.95 & 61.65 & \cellcolor{green!20}$-$19.9\% \\
			1015 & 76.91 & 77.14 & \cellcolor{red!20}+0.3\% \\
			1977 & 76.86 & 82.51 & \cellcolor{red!20}+7.3\% \\
			4000 & 77.05 & 94.29 & \cellcolor{red!20}+22.4\% \\
			\hline
			\textbf{Avg.} & \textbf{76.93} & \textbf{73.14} & \cellcolor{green!20}\textbf{$-$4.9\%} \\
			\hline
		\end{tabular}
	\end{table}
	
	The global energy saving is \textbf{4.9\%} with a maximum of \textbf{62.5\%}
	at 139\,m. The moderate overhead at intermediate and long distances (6--22\%)
	reflects the use of longer-ToA, robust DR configurations (DR8/DR10) required
	to achieve the PDR improvements reported above---a deliberate trade-off
	prioritizing reliability over raw energy consumption.
	
	\subsection{Global Performance Summary}
	
	Table~\ref{tab:summary} consolidates the key gains of our approach.
	
	\begin{table}[!t]
		\centering
		\caption{Summary of DDQN Gains over Static Configurations}
		\label{tab:summary}
		\begin{tabular}{|l|c|c|c|}
			\hline
			\textbf{Metric} & \textbf{Static} & \textbf{DDQN} & \textbf{Gain} \\
			\hline
			Avg.\ PDR (\%)           & 92.7  & 95.0  & \textbf{+2.48\,pts} \\
			PDR at 4000\,m (\%)      & 72.4  & 83.7  & \textbf{+15.6\%} \\
			Avg.\ TX power (dBm)     & 14.0  & 9.5   & \textbf{$-$32\%} \\
			Avg.\ energy/pkt (mJ)    & 76.93 & 73.14 & \textbf{$-$4.9\%} \\
			Min.\ energy/pkt (mJ)    & 76.88 & 28.81 & \textbf{$-$62.5\%} \\
			Favorable distances      & 2/12  & 10/12 & \textbf{83\%} \\
			\hline
		\end{tabular}
	\end{table}
	
	% =====================
	% VI. CONCLUSION
	% =====================
	\section{Conclusion}
	\label{sec:conclusion}
	
	This paper presented the first deep reinforcement learning approach for the
	dynamic adaptation of LR-FHSS transmission parameters. We developed a
	high-fidelity LR-FHSS event-driven simulator (RMSE 12.14\% on PDR,
	6.56\,dB on RSSI), formalized the optimization problem as a Markov Decision
	Process with 76 discrete (DR, power) actions, and trained a DDQN agent on this
	simulator.
	
	Evaluations demonstrate consistent improvements over static configurations:
	+2.48\,pts average PDR (+15.6\% at 4\,km), 32\% transmit power reduction,
	and 4.9\% average energy savings per packet (up to 62.5\% at short range).
	The agent's learned policy---preferring high-throughput DR at short range and
	switching to robust DR at long range, while minimizing transmit power---is
	coherent with communication theory and confirms the relevance of deep
	reinforcement learning for autonomous LR-FHSS optimization.
	
	Future work will address: richer channel models (multi-path fading, dynamic
	shadowing), online learning for real-time policy adaptation, evaluation in
	multi-node dense scenarios, and experimental deployment on physical LR-FHSS
	hardware.
	
	% =====================
	% REFERENCES
	% =====================
	\begin{thebibliography}{99}
		
		\bibitem{Raza2017LPWANSurvey}
		U. Raza, P. Kulkarni, and M. Sooriyabandara,
		``Low Power Wide Area Networks: An Overview,''
		\textit{IEEE Commun. Surveys Tuts.}, vol.~19, no.~2, pp.~855--873, 2017.
		
		\bibitem{loraalliance2020regional}
		LoRa Alliance, ``LoRaWAN Regional Parameters RP002-1.0.0,'' 2020.
		
		\bibitem{loraalliance2025regional}
		LoRa Alliance, ``LoRaWAN Regional Parameters RP002-1.0.5,'' 2025.
		
		\bibitem{ieee_lr_fhss_2021}
		C. Delgado, J. M. Sanz, C. Blondia, and J. Famaey,
		``Towards a New Era of Massive IoT: LR-FHSS for Satellite Networks,''
		in \textit{Proc. IEEE VTC}, 2021.
		
		\bibitem{Li2018ADRanalysis}
		S. Li, U. Raza, and A. Khan,
		``How Agile is the Adaptive Data Rate Mechanism of LoRaWAN?''
		in \textit{Proc. IEEE ICDCS}, 2018, pp.~1360--1365.
		
		\bibitem{greenlora}
		A. Saleem, B. Kantarci, and M. St-Hilaire,
		``GreenLoRa: DDPG-based Adaptive Transmission for LoRaWAN,''
		\textit{IEEE Internet Things J.}, vol.~10, 2023.
		
		\bibitem{greenlora2024}
		A. Saleem et al., ``DP-FH: Deterministic Policy with Frequency Hopping,''
		in \textit{Proc. IEEE WCNC}, 2024.
		
		\bibitem{maldonado2025headerless}
		D. Maldonado et al.,
		``Header-less LR-FHSS Decoding for Dense Networks,''
		\textit{IEEE Commun. Lett.}, 2025.
		
		\bibitem{maldonado2025enhancing}
		D. Maldonado et al.,
		``Enhancing LR-FHSS via FHS Optimization and Demodulator Allocation,''
		\textit{IEEE Trans. Wireless Commun.}, 2025.
		
		\bibitem{ullah_extending_doppler_2025}
		S. Ullah et al.,
		``Extending LR-FHSS to Direct-to-Satellite with Doppler Pre-Compensation,''
		\textit{IEEE Internet Things J.}, 2025.
		
		\bibitem{lrfhss_comparaison}
		V. Sanchez-Vital, G. Chalhoub, and A. Freitas,
		``Experimental Comparison of LoRa and LR-FHSS,''
		in \textit{Proc. IEEE PIMRC}, 2023.
		
		\bibitem{sanchez_vital_2025_data}
		V. Sanchez-Vital, ``LR-FHSS Measurement Dataset,'' Zenodo, 2024.
		
		\bibitem{lr_fhss_overview}
		O. Afisiadis, M. Cotting, A. Burg, and A. Balatsoukas-Stimming,
		``LR-FHSS: Overview and Performance Analysis,''
		\textit{IEEE Commun. Lett.}, vol.~26, no.~2, pp.~264--268, 2022.
		
		\bibitem{arxiv_lrfhss_2020}
		N. Podevijn et al.,
		``Analytical Study of LR-FHSS Capacity,'' \textit{arXiv:2007.02909}, 2020.
		
		\bibitem{lr_fhss_traces}
		R. Eletr et al.,
		``Analysis of LR-FHSS Real Network Traces,''
		in \textit{Proc. EWSN}, 2023.
		
		\bibitem{lfsr}
		R. Lidl and H. Niederreiter,
		\textit{Introduction to Finite Fields and Their Applications}.
		Cambridge Univ. Press, 1994.
		
		\bibitem{semtech2022lrfhss}
		Semtech, ``AN1200.64: LR-FHSS Application Note,'' 2022.
		
		\bibitem{puterman1994mdp}
		M. L. Puterman,
		\textit{Markov Decision Processes}. Wiley, 1994.
		
		\bibitem{sutton2018reinforcement}
		R. S. Sutton and A. G. Barto,
		\textit{Reinforcement Learning: An Introduction}, 2nd~ed. MIT Press, 2018.
		
		\bibitem{mnih2015dqn}
		V. Mnih et al.,
		``Human-level control through deep reinforcement learning,''
		\textit{Nature}, vol.~518, pp.~529--533, 2015.
		
		\bibitem{vanhasselt2016ddqn}
		H. van Hasselt, A. Guez, and D. Silver,
		``Deep Reinforcement Learning with Double Q-Learning,''
		in \textit{Proc. AAAI}, 2016.
		
		\bibitem{ullah2024_energy}
		S. Ullah et al.,
		``Energy Characterization of LR-FHSS IoT Devices,''
		\textit{IEEE Sensors Lett.}, 2024.
		
		\bibitem{SanchezVital2024EnergyLRFHSS}
		V. Sanchez-Vital et al.,
		``Energy Efficiency Analysis of LR-FHSS for Low-Power IoT,''
		\textit{IEEE Internet Things J.}, 2024.
		
		\bibitem{rappaport2002wireless}
		T. S. Rappaport,
		\textit{Wireless Communications: Principles and Practice}, 2nd~ed.
		Prentice Hall, 2002.
		
		\bibitem{nair2010relu}
		V. Nair and G. E. Hinton,
		``Rectified Linear Units Improve Restricted Boltzmann Machines,''
		in \textit{Proc. ICML}, 2010.
		
		\bibitem{srivastava2014dropout}
		N. Srivastava et al.,
		``Dropout: A Simple Way to Prevent Overfitting,''
		\textit{JMLR}, vol.~15, pp.~1929--1958, 2014.
		
		\bibitem{lr_fhss_dts}
		G. Cochet et al.,
		``LR-FHSS for Direct-to-Satellite IoT,''
		in \textit{Proc. IEEE GLOBECOM}, 2022.
		
		\bibitem{free}
		M. Bor, U. Roedig, T. Voigt, and J. Alonso,
		``Do LoRa Low-Power Wide-Area Networks Scale?''
		in \textit{Proc. ACM MSWiM}, 2016.
		
		\bibitem{Semtech_LR1121_Datasheet_2023}
		Semtech, ``LR1121 Datasheet,'' 2023.
		
	\end{thebibliography}
	
\end{document}